\documentclass[11pt,a4paper]{article}

% ---- Packages -------------------------------------------------------------
\usepackage[utf8]{inputenc}
\usepackage[T1]{fontenc}
\usepackage{lmodern}
\usepackage[margin=2.4cm]{geometry}
\usepackage{xcolor}
\usepackage{listings}
\usepackage{fancyhdr}
\usepackage{titlesec}
\usepackage{enumitem}
\usepackage{booktabs}
\usepackage{tabularx}
\usepackage{tcolorbox}
\usepackage{hyperref}
\usepackage{parskip}
\tcbuselibrary{skins,breakable}

% ---- Colors (matches the app's neutral/OKLCH palette) ---------------------
\definecolor{ink}{HTML}{111111}
\definecolor{accent}{HTML}{DF0000}      % Wails red
\definecolor{go}{HTML}{00ADD8}          % Go cyan
\definecolor{react}{HTML}{61DAFB}       % React blue
\definecolor{muted}{HTML}{6B7280}
\definecolor{codebg}{HTML}{F6F7F9}
\definecolor{codeframe}{HTML}{E5E7EB}
\definecolor{kw}{HTML}{7C3AED}
\definecolor{str}{HTML}{047857}
\definecolor{cmt}{HTML}{9CA3AF}

% ---- Hyperref -------------------------------------------------------------
\hypersetup{
  colorlinks=true, linkcolor=accent, urlcolor=go, citecolor=accent,
  pdftitle={Desktop UI Stack Guide}, pdfauthor={rexionmars}
}

% ---- Headings -------------------------------------------------------------
\titleformat{\section}{\Large\bfseries\color{ink}}{\thesection}{0.6em}{}
\titleformat{\subsection}{\large\bfseries\color{ink}}{\thesubsection}{0.6em}{}
\titlespacing*{\section}{0pt}{1.4em}{0.5em}

% ---- Header / footer ------------------------------------------------------
\pagestyle{fancy}\fancyhf{}
\lhead{\small\color{muted}Desktop UI Stack Guide}
\rhead{\small\color{muted}Wails + Go + React}
\cfoot{\small\color{muted}\thepage}
\renewcommand{\headrulewidth}{0.2pt}

% ---- Listings -------------------------------------------------------------
\lstdefinestyle{base}{
  backgroundcolor=\color{codebg},
  basicstyle=\ttfamily\footnotesize\color{ink},
  keywordstyle=\color{kw}\bfseries,
  stringstyle=\color{str},
  commentstyle=\color{cmt}\itshape,
  frame=single, framerule=0.4pt, rulecolor=\color{codeframe},
  breaklines=true, showstringspaces=false,
  numbers=none, tabsize=2, xleftmargin=6pt, xrightmargin=4pt,
  aboveskip=8pt, belowskip=4pt,
}
\lstdefinelanguage{Go}{
  keywords={package,import,func,var,type,struct,interface,return,for,range,if,else,go,defer,chan,map,const,nil,true,false},
  sensitive=true, comment=[l]{//}, morestring=[b]", morestring=[b]`,
}
\lstdefinelanguage{TypeScript}{
  keywords={import,export,from,const,let,var,function,return,await,async,type,interface,as,new},
  sensitive=true, comment=[l]{//}, morecomment=[s]{/*}{*/}, morestring=[b]", morestring=[b]',
}
\lstdefinelanguage{CSS}{
  keywords={import,theme,inline,root,dark},
  sensitive=true, morecomment=[s]{/*}{*/}, morestring=[b]", morestring=[b]',
}
\lstdefinelanguage{bash}{
  keywords={go,install,npx,wails,build,dev,add,init},
  sensitive=true, comment=[l]{\#}, morestring=[b]", morestring=[b]',
}
\lstset{style=base}

% ---- Callout box ----------------------------------------------------------
\newtcolorbox{callout}[1][]{
  colback=codebg, colframe=accent, boxrule=0.8pt, arc=2pt,
  left=8pt, right=8pt, top=6pt, bottom=6pt, breakable, #1
}

% ---- Title data -----------------------------------------------------------
\title{\vspace{-1.2cm}\bfseries\color{ink}\Huge Desktop UI Stack Guide\\[4pt]
  \large\color{muted}\normalfont A reusable blueprint for cross-platform desktop apps\\
  \normalsize with \textbf{Wails\,+\,Go\,+\,React} (the ``Electron alternative'')}
\author{}
\date{\small\color{muted}Reference stack extracted from the TikTok Downloader Desktop project}

% ===========================================================================
\begin{document}
\maketitle
\thispagestyle{fancy}
\vspace{-0.8cm}

\begin{callout}
\textbf{What this is.} A pragmatic, copy-paste-ready description of a modern
desktop UI stack: a single Go binary that renders a React frontend inside the
operating system's native WebView. No Chromium bundled, so binaries are
\textbf{$\sim$6\,MB} instead of $\sim$150\,MB. Use it as the starting template
for new projects.
\end{callout}

% ---------------------------------------------------------------------------
\section{Mental model}

The app is one native executable. The Go side owns the window and the business
logic; the web side (React, built by Vite) is \emph{embedded} into that binary
and drawn by the OS WebView. A generated, type-safe bridge lets the frontend
call Go methods as if they were local async functions.

\begin{center}
\begin{tcolorbox}[colback=white,colframe=codeframe,boxrule=0.6pt,arc=3pt,width=\linewidth]
\ttfamily\small
Single Go binary\\
\hspace*{1em}|\\
\hspace*{1em}+--\,Native WebView\ \ (WebView2 / WKWebView / WebKitGTK)\\
\hspace*{1em}|\hspace*{4em}+--\,React 19 + Vite\ \ (frontend/dist embedded via go:embed)\\
\hspace*{1em}|\\
\hspace*{1em}\textcolor{accent}{$\updownarrow$ generated bridge:}\ window.go.main.App\ \ (typed Promises)\\
\hspace*{1em}|\\
\hspace*{1em}+--\,App struct (Go)\ \ --\ methods exposed to the frontend
\end{tcolorbox}
\end{center}

\begin{callout}
\textbf{Why not Electron?} Electron ships a full Chromium + Node runtime per app.
Wails reuses the WebView already installed on the OS and compiles Go to a native
binary. Result: dramatically smaller size, lower memory, native performance ---
at the cost of small cross-platform WebView differences you must test for.
\end{callout}

% ---------------------------------------------------------------------------
\section{The two layers}

\subsection{Layer 1 --- Desktop shell (Go)}
\begin{center}
\begin{tabularx}{\linewidth}{@{}l l X@{}}
\toprule
\textbf{Piece} & \textbf{Version} & \textbf{Role} \\
\midrule
Wails v2   & \texttt{v2.10.1} & Window, embed frontend, generate Go$\leftrightarrow$JS bridge \\
Go         & \texttt{1.23+}   & Backend + business logic \\
chromedp   & \texttt{v0.11.2} & \emph{Optional} --- browser automation / scraping \\
\bottomrule
\end{tabularx}
\end{center}

\subsection{Layer 2 --- Frontend (web)}
\begin{center}
\begin{tabularx}{\linewidth}{@{}l l X@{}}
\toprule
\textbf{Category} & \textbf{Library / version} & \textbf{Why} \\
\midrule
Framework      & React \texttt{19.2}                 & UI runtime \\
Build tool     & Vite \texttt{7.3}                   & Fast dev server, HMR, bundling \\
React plugin   & @vitejs/plugin-react \texttt{5.1}   & JSX + Fast Refresh \\
Language       & TypeScript \texttt{5.9}             & Type-safety, including the bridge \\
Styling        & Tailwind CSS \texttt{4.1}           & Utility-first, \textbf{no config file} \\
Tailwind/Vite  & @tailwindcss/vite \texttt{4.1}      & Official plugin (replaces PostCSS) \\
Components     & shadcn/ui                            & Components you \emph{own} (copied in) \\
Primitives     & Radix UI \texttt{1.x}               & Headless accessibility (a11y) \\
Icons          & lucide-react \texttt{0.562}         & SVG icon set \\
Toasts         & sonner \texttt{2.0}                 & Notifications \\
Theming        & next-themes \texttt{0.4}            & Light / dark / system \\
CSS animation  & tw-animate-css \texttt{1.4}         & Keyframe utilities \\
Class merge    & clsx + tailwind-merge               & The \texttt{cn()} helper \\
Variants       & class-variance-authority \texttt{0.7} & Component variants (\texttt{cva}) \\
\bottomrule
\end{tabularx}
\end{center}

% ---------------------------------------------------------------------------
\section{Five concepts that make the stack work}

\subsection{1.\ The generated Go$\leftrightarrow$JS bridge}
List your Go struct in \texttt{Bind} and Wails generates typed TypeScript
wrappers. Struct types become TS interfaces. Type-safety crosses the
native$\leftrightarrow$web boundary.

\begin{lstlisting}[language=Go]
// main.go
Bind: []interface{}{ app },
\end{lstlisting}

\begin{lstlisting}[language=TypeScript]
// frontend/wailsjs/go/main/App.d.ts  (generated)
export function DownloadURL(req: main.DownloadRequest): Promise<backend.Result>;

// in a React component
const res = await DownloadURL({ url, output, hd, watermark });
\end{lstlisting}

\begin{callout}
\textbf{Gotcha.} A \texttt{nil} Go slice serializes to JSON \texttt{null} (not
\texttt{[]}), which breaks \texttt{.length}/\texttt{.map} in React. Always
initialize returned slices as \texttt{[]T\{\}} at the Go boundary, and defensively
use \texttt{?? []} on the JS side.
\end{callout}

\subsection{2.\ Tailwind 4 without a config file + OKLCH tokens}
Tailwind 4 drops \texttt{tailwind.config.js}. Everything lives in the CSS entry
file. Design tokens use OKLCH for perceptually uniform, consistent theming.

\begin{lstlisting}[language=CSS]
/* src/index.css */
@import "tailwindcss";
@theme inline { --color-primary: var(--primary); }
:root  { --primary: oklch(0.205 0 0); }   /* light */
.dark  { --primary: oklch(0.922 0 0); }   /* dark  */
\end{lstlisting}

\subsection{3.\ shadcn/ui --- you own the components}
shadcn is \emph{not} an npm dependency. You copy each component's source into
\texttt{src/components/ui/} (\texttt{button.tsx}, \texttt{input.tsx}, \dots).
They are Radix (a11y) + Tailwind + \texttt{cva} (variants). Because the code is
yours, you edit it freely. Configuration lives in \texttt{components.json}.

\subsection{4.\ The \texttt{cn()} helper}
The glue of the whole stack: merge conditional classes and resolve Tailwind
conflicts.

\begin{lstlisting}[language=TypeScript]
// src/lib/utils.ts
import { clsx, type ClassValue } from "clsx";
import { twMerge } from "tailwind-merge";
export function cn(...inputs: ClassValue[]) { return twMerge(clsx(inputs)); }
\end{lstlisting}

\subsection{5.\ WebView app patterns (non-obvious)}
\begin{itemize}[leftmargin=1.2em,itemsep=2pt]
  \item \textbf{Frameless window + custom title bar}: \texttt{Frameless: true}
        in Go plus a \texttt{<TitleBar>} with a draggable region via CSS
        (\texttt{-{}-wails-draggable: drag}).
  \item \textbf{External links}: intercept and open in the OS browser
        (\texttt{BrowserOpenURL}); otherwise they open \emph{inside} the app.
  \item \textbf{Go$\to$JS push} for progress: \texttt{runtime.EventsEmit} in Go
        $\leftrightarrow$ \texttt{EventsOn} in React. Methods = request/response;
        events = streaming.
  \item \textbf{External images} in the WebView may need
        \texttt{referrerPolicy="no-referrer"} (hotlink / CORS).
\end{itemize}

% ---------------------------------------------------------------------------
\section{Project structure (template)}

\begin{lstlisting}
project/
|-- main.go              # window, go:embed, Bind
|-- app.go               # App struct -> methods exposed to the frontend
|-- backend/             # Go business logic
|-- wails.json           # name, frontend build commands
`-- frontend/
    |-- vite.config.ts   # @/ alias, injects version from wails.json
    |-- components.json   # shadcn config
    |-- src/
    |   |-- index.css     # Tailwind + OKLCH tokens
    |   |-- main.tsx      # ThemeProvider + Toaster
    |   |-- App.tsx       # layout: TitleBar + Sidebar + pages
    |   |-- components/ui/ # shadcn components (you own)
    |   |-- components/    # your feature components
    |   `-- lib/utils.ts   # cn()
    `-- wailsjs/          # generated Go<->JS bridge
\end{lstlisting}

% ---------------------------------------------------------------------------
\section{Commands}

\begin{lstlisting}[language=bash]
# one-time setup
go install github.com/wailsapp/wails/v2/cmd/wails@latest
npx shadcn@latest init
npx shadcn@latest add button input switch progress

# development (hot-reload: Vite watches frontend, Go rebuilds)
wails dev

# native build
wails build -platform darwin/universal   # .app
wails build -platform windows/amd64      # .exe
wails build -platform linux/amd64        # binary
\end{lstlisting}

% ---------------------------------------------------------------------------
\section{Pitfalls \& trade-offs}

\begin{enumerate}[leftmargin=1.4em,itemsep=3pt]
  \item \textbf{The ``god component''.} \texttt{App.tsx} tends to accumulate all
        state and orchestration. Extract logic into hooks early
        (\texttt{useX}) and keep components presentational.
  \item \textbf{WebView $\neq$ identical Chrome} across OSes. Test on all three
        targets; WKWebView (macOS) sometimes differs subtly.
  \item \textbf{macOS signing.} Unsigned \texttt{.app}/\texttt{.dmg} triggers the
        ``unidentified developer'' warning; notarization needs a paid Apple
        Developer account (\$99/yr).
  \item \textbf{External images} may require \texttt{referrerPolicy} due to
        hotlink/CORS rules.
  \item \textbf{CI matrix.} Build on all three platforms via GitHub Actions and
        publish binaries on \texttt{v*} tags (a \texttt{build.yml} matrix is the
        canonical template).
\end{enumerate}

\vfill
\begin{center}\small\color{muted}
Blueprint distilled from a working project --- verified to build and run on macOS,
Windows and Linux.
\end{center}

\end{document}
